\documentclass[12pt]{article}
\usepackage{amsmath,amssymb,amsfonts}
\usepackage[margin=1in]{geometry}

\begin{document}

\begin{center}
{\large\bfseries 1st Irish Mathematical Olympiad\\
30 April 1988}
\end{center}

\bigskip
\noindent\textbf{Paper 1}
\bigskip

\begin{enumerate}
\item A pyramid with a square base, and all its edges of length 2, is joined to a regular tetrahedron, whose edges are also of length 2, by gluing together two of the triangular faces. Find the sum of the lengths of the edges of the resulting solid.

\item $A,B,C,D$ are the vertices of a square, and $P$ is a point on the arc $CD$ of its circumcircle. Prove that \[|PA|^2 - |PB|^2 = |PB|\cdot|PD| - |PA|\cdot|PC|.\]

\item $ABC$ is a triangle inscribed in a circle, and $E$ is the mid-point of the arc subtended by $BC$ on the side remote from $A$. If through $E$ a diameter $ED$ is drawn, show that the measure of the angle $DEA$ is half the magnitude of the difference of the measures of the angles at $B$ and $C$.

\item A mathematical moron is given the values $b, c, A$ for a triangle $ABC$ and is required to find $a$. He does this by using the cosine rule \[a^2 = b^2 + c^2 - 2bc\cos A\] and misapplying the law of the logarithm to this to get \[\log a^2 = \log b^2 + \log c^2 - \log(2bc\cos A).\] He proceeds to evaluate the right-hand side correctly, takes the anti-logarithms and gets the correct answer. What can be said about the triangle $ABC$?

\item A person has seven friends and invites a different subset of three friends to dinner every night for one week (seven days). In how many ways can this be done so that all friends are invited at least once?

\item Suppose you are given $n$ blocks, each of which weighs an integral number of pounds, but less than $n$ pounds. Suppose also that the total weight of the $n$ blocks is less than $2n$ pounds. Prove that the blocks can be divided into two groups, one of which weighs exactly $n$ pounds.

\item A function $f$, defined on the set of real numbers $\mathbb{R}$ is said to have a \emph{horizontal chord} of length $a > 0$ if there is a real number $x$ such that $f(a + x) = f(x)$. Show that the cubic \[f(x) = x^3 - x \qquad (x \in \mathbb{R})\] has a horizontal chord of length $a$ if, and only if, $0 < a \le 2$.

\item Let $x_1, x_2, x_3, \ldots$ be a sequence of nonzero real numbers satisfying \[x_n = \frac{x_{n-2}x_{n-1}}{2x_{n-2} - x_{n-1}}, \qquad n = 3, 4, 5, \ldots\] Establish necessary and sufficient conditions on $x_1, x_2$ for $x_n$ to be an integer for infinitely many values of $n$.

\item The year 1978 was ``peculiar'' in that the sum of the numbers formed with the first two digits and the last two digits is equal to the number formed with the middle two digits, i.e., $19 + 78 = 97$. What was the last previous peculiar year, and when will the next one occur?

\item Let $0 \le x \le 1$. Show that if $n$ is any positive integer, then \[(1+x)^n \ge (1-x)^n + 2nx(1-x^2)^{\frac{n-1}{2}}.\]

\item If facilities for division are not available, it is sometimes convenient in determining the decimal expansion of $1/a$, $a > 0$, to use the iteration \[x_{k+1} = x_k(2 - ax_k), \qquad k = 0, 1, 2, \ldots,\] where $x_0$ is a selected ``starting'' value. Find the limitations, if any, on the starting values $x_0$, in order that the above iteration converges to the desired value $1/a$.

\item Prove that if $n$ is a positive integer, then \[\sum_{k=1}^{n} \cos^4\left(\frac{k\pi}{2n+1}\right) = \frac{6n-5}{16}.\]

\end{enumerate}

\newpage

\begin{center}
{\large\bfseries 1st Irish Mathematical Olympiad\\
30 April 1988}
\end{center}

\bigskip
\noindent\textbf{Paper 2}
\bigskip

\begin{enumerate}
\setcounter{enumi}{12}
\item The triangles $ABG$ and $AEF$ are in the same plane. Between them the following conditions hold: \begin{enumerate} \item[(a)] $E$ is the mid-point of $AB$; \item[(b)] points $A, G$ and $F$ are on the same line; \item[(c)] there is a point $C$ at which $BG$ and $EF$ intersect; \item[(d)] $|CE| = 1$ and $|AC| = |AE| = |FG|$. \end{enumerate} Show that if $|AG| = x$, then $|AB| = x^3$.

\item Let $x_1, \ldots, x_n$ be $n$ integers, and let $p$ be a positive integer, with $p < n$. Put \begin{align*} S_1 &= x_1 + x_2 + \ldots + x_p, \\ T_1 &= x_{p+1} + x_{p+2} + \ldots + x_n, \\ S_2 &= x_2 + x_3 + \ldots + x_{p+1}, \\ T_2 &= x_{p+2} + x_{p+3} + \ldots + x_n + x_1, \\ &\vdots \\ S_n &= x_n + x_1 + x_2 + \ldots + x_{p-1}, \\ T_n &= x_p + x_{p+1} + \ldots + x_{n-1}. \end{align*} For $a = 0,1,2,3$, and $b = 0,1,2,3$, let $m(a,b)$ be the number of numbers $i$, $1 \le i \le n$, such that $S_i$ leaves remainder $a$ on division by 4 and $T_i$ leaves remainder $b$ on division by 4. Show that $m(1,3)$ and $m(3,1)$ leave the same remainder when divided by 4 if, and only if, $m(2,2)$ is even.

\item A city has a system of bus routes laid out in such a way that \begin{enumerate} \item[(a)] there are exactly 11 bus stops on each route; \item[(b)] it is possible to travel between any two bus stops without changing routes; \item[(c)] any two bus routes have exactly one bus stop in common. \end{enumerate} What is the number of bus routes in the city?

\end{enumerate}

\end{document}
